\section{伪部分子分布}

准部分子分布 $Q(y,P)$ 的复杂结构，形式上源于它是由函数
${\cal M} (\nu, \nu^2/P^2)$ 对 $\nu$ 作 Fourier 变换得到的——在该函数中，
$\nu$ 同时出现在 Ioffe 时间分布的第一与第二宗量中。
因此必须取足够大的 $P$ 值，
才能忽略来自第二宗量的 $\nu$ 依赖。

另一条路 \cite{Radyushkin:2017cyf} 是设法{\it 消除}
由 ${\cal M} (\nu, z_3^2)$ 带来的 \mbox{$z_3^2$ 依赖}。其主要想法基于如下观察：
若取修正后的函数 ${\cal M} (\nu, z_3^2)/D(z_3^2)$ 的 $\nu$-Fourier 变换，
则只要 $D(z_3^2)$ 仅是 $z_3^2$ 的函数（不含 $\nu$）且满足
\mbox{$z_3^2=0$ 时等于 1}，其 $z_3\to 0$ 极限给出的 PDF 就与原始 Ioffe 时间分布相同。

于是策略便是：寻找一个函数 $D(z_3^2)$，使其 \mbox{$z_3^2$ 依赖}
能在最大程度上抵消 ${\cal M} (\nu, z_3^2)$ 的 $z_3^2$ 依赖；
下一步是用 $z_3^2$ 的多项式拟合残余的多项式 $z_3^2$ 依赖
（对不同的 $\nu$ 可以采用不同的多项式），
从而把数据外推到 $z_3^2=0$ 极限。
所得函数的 Fourier 变换便对应于
原始 Ioffe 时间分布 ${\cal M} (\nu, z_3^2)$ 在 $z_3^2 \to 0$ 极限下的同一 PDF。


在最理想的情况下，比值 ${\cal M} (\nu, z_3^2)/D(z_3^2)$
不含多项式型 $z_3^2$ 依赖（或仅有非常轻微的依赖）。
特别地，当 ${\cal M} (\nu, z_3^2)$ 因子化，即
 ${\cal M} (\nu, z_3^2)= {\cal M} (\nu, 0) {\cal M} (0, z_3^2)$ 时，
应取 \mbox{$D(z_3^2) ={\cal M} (0, z_3^2)$}。此时约化函数
\begin{align}
{\mathfrak M} (\nu, z_3^2) \equiv \frac{ {\cal M} (\nu, z_3^2)}{{\cal M} (0, z_3^2)} \
 \label{redm}
\end{align}
就等于 $ {\cal M} (\nu, 0)$，
求取 $z_3\to 0$ 极限的任务即告完成。

虽然这样的因子化并没有“第一性原理”层面的依据，
但可以预期：不同 $\nu$ 的函数 ${\cal M} (\nu, z_3^2)$
对 $z_3$ 的依赖大同小异，
这基本上反映的是核子的有限尺寸。

如前所述，只要 TMD ${\cal F} (x,k_\perp^2)$ 的软部分因子化，
${\cal M} (\nu, z_3^2)$ 的软部分就会因子化。
这是 TMD 从业者的标准假设
（例如见文献 \cite{Anselmino:2013lza}）。
因此 ${\cal M} (\nu, z_3^2)$ 的这部分 $z_3^2$ 依赖
很有可能被静止系函数 ${\cal M} (0, z_3^2) $ 抵消或大幅削弱。


在格点上还存在另一个（也是麻烦的，如文献 \cite{Ishikawa:2016znu} 所述）
$z_3$ 依赖来源：
由规范链接 ${ \hat E} (0,z_3; A)$ 重正化生成的
因子 $Z(z_3^2)$。
幸运的是，这个麻烦的因子 $Z(z_3^2)$ 不依赖于 $\nu$，
并且对比值 ${\mathfrak M} (\nu, z_3^2) $ 的分子与分母是同一个。
这为使用 ${\cal M} (0, z_3^2)$ 作为因子 $D(z_3^2)$ 提供了又一条理由。

 于是，该方案建议对约化 Ioffe 时间函数 ${\mathfrak M} (\nu, z_3^2)$ 进行格点研究。
 即便它带有残余的多项式 $z_3^2$ 依赖，
 把这一依赖外推到 $z_3=0$ 也应当容易得多，
 远易于处理原始 Ioffe 时间分布 $ {\cal M} (\nu, z_3^2)$ 的 $z_3$ 依赖。


此外，如果观察到比值 $ {\mathfrak M} (\nu, z_3^2)$ 不含 $z_3$ 依赖，
就应得出结论：$ {\cal  M} (\nu, z_3^2)$ 发生了因子化。
事实上，早在多年以前格点 QCD 横动量分布的开创性研究
\cite{Musch:2010ka} 中就已观察到这种因子化。


不过，仍有一个无法避免的因子化破坏来源。
当 $z_3$ 较小时，
$ {\cal M} (\nu, z_3^2)$ 具有 $\ln z_3^2$ 型对数奇异性，
它们生成了 PDF 的微扰演化。
如前所述，此时 $1/z_3$
类比于标准 OPE 方法中标度依赖的 PDF $f(x,\mu^2)$ 的重正化参数 $\mu$。

    \begin{figure}[t]
   \centerline{ \includegraphics[width=4.0in]{images/MRER.pdf}  }
    \caption{模型分布 ${\cal M} (\nu)$ 的实部及其演化的控制函数 $-B \otimes  {\rm Re} \,  {\cal M}$
    （此处的负号只是为了把两条曲线画在同一张图上）。
    \label{MBM}}
    \end{figure}

更具体地说，对小 $z_3$ 取值，伪部分子分布 ${\cal P} (x,z_3^2)$
满足一个关于 $1/z_3$ 的领头阶演化方程，
它与 $f (x,\mu^2)$ 关于 $\mu$ 的演化方程完全相同。
约化 Ioffe 时间分布 ${\mathfrak M} (\nu, z_3^2)$ 的演化方程也可以写为
\cite{Radyushkin:2017cyf}
    \begin{align}
    \frac{d}{d \ln z_3^2} \,
{\mathfrak M} (\nu, z_3^2)    &= - \frac{\alpha_s}{2\pi} \, C_F
\int_0^1  du \,   B ( u )   {\mathfrak M} (u \nu, z_3^2)
 ,
\label{EE}
 \end{align}
其中 $C_F=4/3$；非单态夸克情形的领头阶演化核 $B(u)$ 为 \cite{Braun:1994jq}
   \begin{align}
B (u)    &=  \left [ \frac{1+u^2}{1- u} \right ]_+  \
 ,
\label{Bu}
 \end{align}
这里 $[ \ldots ]_+$ 表示“加”（plus）规定，即
    \begin{align}
& \int_0^1  du \,    \left [ \frac{1+u^2}{1- u} \right ]_+     {\mathfrak M} (u \nu)
\nn & =
\int_0^1  du \,  \frac{1+u^2}{1- u} \,    [{\mathfrak M} ( \nu)-{\mathfrak  M} (u \nu)  ]\
 .
\label{plus}
 \end{align}


注意，作为 Fourier 变换，
  \begin{align}
  {\cal M} (\nu)  =
   \int_{-1}^1 dx \, f(x) \,
e^{ix\nu} \,  \  ,
 \label{Mf}
\end{align}
即使函数 $f(x)$ 为实（部分子分布正是如此），
Ioffe 时间分布也同时具有实部与虚部。特别地，
   \begin{align}
  {\rm Re} \, {\cal M} (\nu)  =
   \int_{-1}^1 dx \, f(x) \,
\cos (x\nu)  \,  \  ,
 \label{ReMf}
\end{align}
以及
    \begin{align}
  {\rm Im} \, {\cal M} (\nu)  =
   \int_{-1}^1 dx \, f(x) \,
\sin (x\nu)  \,  \  .
 \label{ImMf}
\end{align}



在 \mbox{图~\ref{MBM}} 中，我们对模型 PDF
      \begin{align}
q (x)= \frac{315}{32}  \sqrt{x} (1-x)^3 \theta ( 0\leq x \leq 1) \
\end{align}
画出了函数 ${\rm Re}  \, {\cal M} (\nu)$。
它的积分归一化为 1，且仅对正 $x$ 非零，
对应于无反夸克的情形。
我们将看到，这一特定形式恰好出现在实际格点数据的描述中。
在 \mbox{图~\ref{MBMI}} 中，我们对同一模型 PDF 画出函数
${\rm Im}  \, {\cal M} (\nu)$。


    \begin{figure}[t]
     \centerline{ \includegraphics[width=3.8in]{images/MIEIkey.pdf}}
    \caption{模型 Ioffe 时间分布 $  {\cal M} (\nu)$ 的虚部及其演化的控制函数 $B \otimes  {\rm Im} \,  {\cal M}$。
    \label{MBMI}}
    \end{figure}



我们在这些图中还画出了支配演化的卷积积分，即
\mbox{$-B \otimes {\rm Re}\, {\cal M}(\nu)$}
与 \mbox{$B \otimes {\rm Im}\, {\cal M}(\nu)$}。
读者会注意到，由于矢量流守恒，
$B \otimes {\rm }\,  {\cal M}(\nu)$ 在 $\nu=0$ 处为零。
其后果是：微扰演化不影响静止系密度 ${\cal M} (0,z_3^2)$（它总是实的）。
换言之，$\ln z_3^2$ 项只出现在比值 $ {\mathfrak M} (\nu, z_3^2)$ 的分子
${\cal M} ( \nu, z_3^2)  $ 中，
而不出现在分母 ${\cal M} ( 0, z_3^2)  $ 中。

还要注意，实部的演化总是使 ${{\rm Re} \, \cal M} ( \nu, z_3^2)$
随 $z_3^2$ 增大而{\it 减小}。
虚部的演化样式则更为复杂：
当 \mbox{$\nu \lesssim 5.5$} 时，函数 ${{\rm Im} \, \cal M} ( \nu, z_3^2)$
随 $z_3^2$ 增大而{\it 增大}；
只有当 $\nu \sim 5.5$ 以上时，
演化才导致 ${{\rm Im} \, \cal M} (u \nu, z_3^2)$ 随 $z_3^2$ 减小，
演化样式变得与实部类似。
